\section{پیش‌پردازش داده‌ها}

پیش از به‌کارگیری هر الگوریتم یادگیری ماشین، ضروری است که مجموعه‌داده از نظر پاکیزگی، سازگاری
و مقیاس‌بندی مناسب بررسی و آماده‌سازی شود. در این پروژه، فرایند پیش‌پردازش شامل سه مرحلهٔ اصلی است:

\begin{enumerate}
	\item رسیدگی به مقادیر گمشده،
	\item نرمال‌سازی ویژگی‌های عددی،
	\item تقسیم داده‌ها به مجموعه‌های آموزش و آزمون.
\end{enumerate}

تمامی مراحل پیش‌پردازش با استفاده از کتابخانه‌های \lr{pandas} و \lr{numpy} در محیط
\lr{Jupyter Notebook} پیاده‌سازی شده‌اند. به منظور سهولت در نگه‌داری و استفادهٔ مجدد از کد،
هر بخش در قالب توابع جداگانه و ماژولار تعریف شده است.

\subsection{رسیدگی به مقادیر گمشده}

پس از بارگذاری مجموعه‌دادهٔ تجمیع‌شدهٔ بیماری قلبی از فایل \lr{CSV}، ابتدا با استفاده از
فرمان‌های \lr{df.info()} و \lr{df.isna().sum()} وضعیت مقادیر گمشده و ناسازگاری‌های احتمالی
بررسی شد. برای تضمین کامل بودن داده‌ها و جلوگیری از بروز خطا در مراحل بعدی، یک راهبرد
جایگزینی ساده اما مؤثر اعمال گردید:

\begin{itemize}
	\item برای ویژگی‌های \textbf{عددی} مانند \lr{age}، \lr{trestbps}، \lr{chol}، \lr{thalach} و \lr{oldpeak}
	مقادیر گمشده با \emph{میانه}‌ی ستون جایگزین شدند.
	
	\item برای ویژگی‌های \textbf{دسته‌ای یا دودویی} مانند \lr{sex}، \lr{cp}، \lr{fbs}، \lr{restecg}،
	\lr{exang}، \lr{slope}، \lr{ca} و \lr{thal} مقادیر گمشده با \emph{مد} (پرمقدارترین کلاس) جایگزین شدند.
\end{itemize}

استفاده از میانه برای داده‌های عددی باعث کاهش حساسیت نسبت به مقادیر پرت می‌شود و برای داده‌های
پزشکی که ممکن است شامل مقادیر غیرعادی باشند گزینه‌ای مناسب است. همچنین، به‌کارگیری مد برای
ویژگی‌های دسته‌ای باعث حفظ فراوان‌ترین مقدار و جلوگیری از ایجاد کلاس‌های غیرواقعی می‌شود.

این عملیات در قالب تابع \lr{impute\_missing\_values()} پیاده‌سازی شده است. پس از اجرای آن،
بررسی مجدد داده‌ها نشان داد که هیچ مقدار گمشده‌ای باقی نمانده است.

\subsection{نرمال‌سازی ویژگی‌های عددی}

از آن‌جا که الگوریتم \lr{KNN} مبتنی بر فاصله است، بسیار مهم است که ویژگی‌های عددی در یک مقیاس
مشابه قرار داشته باشند. در غیر این صورت، ویژگی‌هایی که دامنهٔ بزرگ‌تری دارند ممکن است بر
محاسبات فاصله غالب شوند و منجر به نتایج نادرست گردند.

در این پروژه از روش \textbf{مقیاس‌گذاری مین–مکس} برای نرمال‌سازی ویژگی‌های عددی استفاده شده است.
برای یک ویژگی عددی $x$ مقدار نرمال‌شدهٔ آن به صورت زیر محاسبه می‌شود:

\begin{equation}
	x' = \frac{x - x_{\min}}{x_{\max} - x_{\min}},
\end{equation}

که $x_{\min}$ و $x_{\max}$ به ترتیب کوچک‌ترین و بزرگ‌ترین مقدار آن ویژگی در \textbf{مجموعهٔ آموزش}
هستند.

برای جلوگیری از نشت اطلاعات (\lr{data leakage})، مقادیر حداقلی و حداکثری تنها از داده‌های
آموزشی به دست آمده و همان مقادیر برای تبدیل داده‌های آزمون استفاده شده‌اند. این فرایند در تابع
\lr{min\_max\_scale\_train\_test()} پیاده‌سازی شده است که خروجی آن نسخهٔ نرمال‌شدهٔ داده‌های آموزش و
آزمون است.

در صورتی که ستونی عددی مقدار ثابتی داشته باشد (یعنی $x_{\max} = x_{\min}$)، به منظور جلوگیری از تقسیم
بر صفر، مقدار نرمال‌شدهٔ آن ستون برابر صفر قرار داده شده است.

\subsection{تقسیم داده‌ها به آموزش و آزمون}

برای ارزیابی توان تعمیم مدل، مجموعه‌داده به دو بخش آموزش و آزمون تقسیم شد. ابتدا داده‌ها به‌طور
تصادفی به‌هم‌ریخته شدند و سپس با نسبت \textbf{۸۰\% برای آموزش} و \textbf{۲۰\% برای آزمون} تفکیک شدند:

\begin{itemize}
	\item ۸۰ درصد داده‌ها برای آموزش مدل استفاده شدند،
	\item ۲۰ درصد داده‌ها برای آزمون نهایی کنار گذاشته شدند.
\end{itemize}

این تقسیم‌بندی توسط تابع \lr{train\_test\_split\_df()} انجام شد که ورودی آن مجموعه‌دادهٔ کامل
و درصد داده‌های آزمون است. برای تضمین تکرارپذیری (\lr{reproducibility})، یک مقدار بذر ثابت
(\lr{random seed}) تنظیم شد.

پس از انجام نرمال‌سازی و تقسیم داده‌ها، ماتریس‌های \lr{X\_train}، \lr{y\_train}،
\lr{X\_test} و \lr{y\_test} حاصل شدند که در مراحل بعدی برای آموزش و ارزیابی طبقه‌بند \lr{KNN}
مورد استفاده قرار گرفتند.